\subsection{Isotypic components}

DEFINITION 9. --- Let G be a topological group and let \(\pi\) be a continuous irreducible representation of G. Let \(\varrho\) be a continuous representation of G in a separated locally convex space E. The \(\pi\)-isotypic component of \(\varrho\) is called the closure in E of the sum of the spaces of all subrepresentations of \(\varrho\) isomorphic to \(\pi\) . This subspace is denoted \(\mathrm{M}_{\pi}(\varrho)\).

The space \(\mathrm{M}_{\pi}(\varrho)\) is a closed subspace of E, which defines a subrepresentation of \(\varrho\) . This space depends only on the class of \(\pi\) in \(\widehat{G}\) .

PROPOSITION 8. — Let G be a topological group. Let \(\varrho\) be a unitary representation of G in a complex Hilbert space E.

\begin{enumerate}
\def\labelenumi{\alph{enumi})}
\item
  Let \(\pi_1\) and \(\pi_2\) be non-isomorphic irreducible unitary representations of G. The subspaces \(\mathrm{M}_{\pi_1}(\varrho)\) and \(\mathrm{M}_{\pi_2}(\varrho)\) are orthogonal;
\item
  Let \(\varrho'\) be a unitary representation of G and \(u\) a G-morphism from \(\varrho\) to \(\varrho'\) . For every irreducible unitary representation \(\pi\) of G, we have \(u(\mathrm{M}_{\pi}(\varrho)) \subset \mathrm{M}_{\pi}(\varrho')\) .
\end{enumerate}

Let \(E_{1}\) and \(E_{2}\) be subspaces of E defining subrepresentations isomorphic to \(\pi_{1}\) and \(\pi_{2}\) respectively. The orthogonal projector of E with image \(E_{2}\) defines by passing to subspaces an element of \(\mathrm{Hom}_{\mathrm{G}}(\pi_{1},\pi_{2})\) , which is zero (cor. 2 of V, p. 387), which proves that \(E_{2}\) is orthogonal to \(E_{1}\) . This proves assertion \(a\) ).

For assertion \(b\) ), note that \(\mathrm{M}_{\pi}(\varrho)\) is, by definition, the closure in E of the space generated by the elements \(x \in \mathrm{E}\) belonging to a closed subspace \(\mathrm{F} \subset \mathrm{E}\) stable under \(\varrho\) such that the subrepresentation \(\varrho_{\mathrm{F}}\) of \(\varrho\) in F is isomorphic to \(\pi\) . It therefore suffices to prove that, in this case, one has \(u(x) \in \mathrm{M}_{\pi}(\varrho')\) . Let \(\mathrm{H} = u(\mathrm{F})\) ; it is a closed subspace of the space of \(\varrho'\) (cor. 5 of V, p. 388) stable under \(\varrho'\) , and u, induced by passage to the subspaces, is a surjective G-morphism from F to H. If H is not zero, then this G-morphism is an isomorphism by Corollary 2 of V, p. 387. The representation of G on H is therefore isomorphic to \(\pi\) , whence it follows that \(u(x)\) belongs to \(\mathrm{M}_{\pi}(\varrho')\) .

For every vector space H and every family \((\mathrm{H}_{i})_{i\in\mathrm{I}}\) of subspaces of H, we say that the spaces \((\mathrm{H}_{i})_{i\in\mathrm{I}}\) are in direct sum if the family \((\mathrm{H}_{i})_{i\in\mathrm{I}}\) satisfies the equivalent conditions of Prop. 11 of A, II, p. 18.

PROPOSITION 9. --- Let G be a topological group and let \(\varrho\) be a unitary representation of G in a complex Hilbert space E. Let \(\pi\) be an irreducible unitary representation of G in a complex Hilbert space \(\mathrm{E}_{\pi}\) . Denote by \(v\) the linear map from \(\mathrm{Hom}_{\mathrm{G}}(\pi, \varrho) \otimes \mathrm{E}_{\pi}\) into E such that \(v(u \otimes x) = u(x)\) for all \((u, x) \in \mathrm{Hom}_{\mathrm{G}}(\pi, \varrho) \times \mathrm{E}_{\pi}\) .

The linear map \(v\) is injective and its image is the sum of the spaces of all subrepresentations of \(\varrho\) which are isomorphic to \(\pi\) . In particular, the image of \(v\) is dense in \(M_{\pi}(\varrho)\) .

Corollary 5 of V, p. 388 implies that the image of v is the sum of the spaces of all subrepresentations of \(\varrho\) which are isomorphic to \(\pi\) .

Let us prove that \(v\) is injective. Let \((u_i)_{i \in I}\) be a basis of the vector space Hom \(_G(\pi, \varrho)\) . For \(i \in I\) , denote by F \(_i\) the image of \(u_i\) and \(\widetilde{u}_i \colon E_\pi \to F_i\) the G-isomorphism deduced from \(u_i\) by passage to the subspaces.

Let us first prove, by induction on the cardinality of a finite subset J of I, that the subspaces \((\mathrm{F}_{i})_{i\in\mathrm{J}}\) are in direct sum.

The case where J is empty is immediate. Suppose that J is nonempty and that the required property holds for subsets of I of cardinality at most \(\text{Card(J)} - 1\) .

Let \(j\) be a fixed element of J. The induction hypothesis implies that the subspaces \(\mathrm{F}_{i}\) for \(i\in\mathrm{J}-\{j\}\) are in direct sum. Let \(\mathrm{F}^{\prime}\) be their sum; it now suffices to prove that \(\mathrm{F}_{j}\cap\mathrm{F}^{\prime}=\{0\}\) .

Suppose that the intersection of \(F_{j}\) and of \(F'\) is not reduced to 0; this intersection is a subrepresentation of \(F_{j}\) , and since the latter is irreducible, we therefore have \(F_{j} \cap F' = F_{j}\) , whence \(F_{j} \subset F'\) .

For \(i \in \mathrm{J} - \{j\}\) , denote by \(\operatorname{pr}_i: \mathrm{F}' \to \mathrm{F}_i\) the projection and by \(\iota_i: \mathrm{F}_i \to \mathrm{F}'\) the inclusion. We have a canonical isomorphism

\[
\operatorname{Hom} _ {\mathrm{G}} \left(\mathrm{F} _ {j}, \mathrm{F} ^ {\prime}\right)\rightarrow \bigoplus_ {i \in \mathrm{J} - \{j \}} \operatorname{Hom} _ {\mathrm{G}} \left(\mathrm{F} _ {j}, \mathrm{F} _ {i}\right)
\]

(formula (1), p. 377) such that \(u \colon \mathrm{F}_j \to \mathrm{F}'\) has as image the family \((\mathrm{pr}_i \circ u)_{i \in \mathrm{J} - \{j\}}\) (A, II, p. 13, cor. 1). Corollary 2 of V, p. 387 implies that \(\widetilde{u}_i \circ \widetilde{u}_j^{-1}\) , which is nonzero, is a basis of the space \(\mathrm{Hom}_{\mathrm{G}}(\mathrm{F}_j, \mathrm{F}_i)\) .

Therefore, the family of G-morphisms \((\iota_{i} \circ \widetilde{u}_{i} \circ \widetilde{u}_{j}^{-1})_{i \in \mathrm{J} - \{j\}}\) is a basis of \(\mathrm{Hom}_{\mathrm{G}}(\mathrm{F}_{j}, \mathrm{F}')\) .

Denote by \(\iota\) the inclusion of \(F_{j}\) in \(F'\) ; it is an element of \(\mathrm{Hom}_{\mathrm{G}}(\mathrm{F}_{j}, \mathrm{F}')\) , hence it is a linear combination of the G-morphisms \(\iota_{i} \circ \widetilde{u}_{i} \circ \widetilde{u}_{j}^{-1}\) for \(i \in \mathrm{J} - \{j\}\) . It follows that \(u_{j}\) is a linear combination of the maps \(u_{i}\) for \(i \neq j\) , which contradicts the linear independence of the family \((u_{i})_{i \in \mathrm{I}}\) . The assertion is therefore proved by induction.

Let us now prove that v is injective. Let w be an element of \(\mathrm{Hom}_{\mathrm{G}}(\pi,\varrho)\otimes\mathrm{E}_{\pi}\) . There exists a unique family \((x_{i})_{i\in\mathrm{I}}\) in \(E_{\pi}\) with finite support such that

\[
w = \sum_ {i \in \mathrm{I}} u _ {i} \otimes x _ {i}
\]

(A, II, p. 62, cor. 1) and we then have

\[
v (w) = \sum_ {i \in I} u _ {i} (x _ {i}).
\]

From what precedes, the condition \(v(w) = 0\) therefore implies that \(u_{i}(x_{i}) = 0\) for all \(i \in I\) , whence \(x_{i} = 0\) for all i since \(u_{i}\) is injective, and therefore w = 0.

PROPOSITION 10. --- Let G be a topological group. Let \(\pi\) be an irreducible unitary representation of G in a complex Hilbert space E \(_{\pi}\) and let \(\varrho\) be a unitary representation of G in a complex Hilbert space E.

There exist families \((\mathrm{E}_{i})_{i\in\mathrm{I}}\) of closed invariant subspaces of E such that the subrepresentation of \(\varrho\) in \(E_{i}\) is isomorphic to \(\pi\) for all \(i\in I\) and such that \(\mathrm{M}_{\pi}(\varrho)\) is the Hilbert sum of the spaces \(E_{i}\) . Moreover, the cardinality of I is independent of the family \((\mathrm{E}_{i})_{i\in\mathrm{I}}\) satisfying these properties.

Let us prove the existence of families satisfying the stated conditions. Let \(\mathcal{O}\) be the set of subsets C of \(\mathrm{Hom}_{\mathrm{G}}(\pi, \varrho)\setminus\{0\}\) such that the subspaces \(u(\mathrm{E}_{\pi})\) for \(u \in \mathrm{C}\) are pairwise orthogonal. The set \(\mathcal{O}\) is ordered by inclusion. It is nonempty because the empty set is a member, and it is of finite character (E, III, p. 34, def. 2) since C belongs to \(\mathcal{O}\) if and only if every two-element subset of C belongs to \(\mathcal{O}\) . By E, III, p. 35, th. 1, there exists a maximal element C of \(\mathcal{O}\) . Let F be the closure of the subspace generated by the spaces \(u(\mathrm{E}_{\pi})\) for \(u \in \mathrm{C}\) ; it is the Hilbert sum of the spaces \(u(\mathrm{E}_{\pi})\) for \(u \in C\) . We shall prove that \(\mathrm{F} = \mathrm{M}_{\pi}(\varrho)\) , which will prove that the family \((u(\mathrm{E}_{\pi}))_{u \in \mathrm{C}}\) has the required properties.

By definition, \(u(\mathrm{E}_{\pi}) \subset \mathrm{M}_{\pi}(\varrho)\) for all \(u \in \mathrm{C}\) , hence \(\mathrm{F}\) is contained in \(\mathrm{M}_{\pi}(\varrho)\) . To prove the reverse inclusion, it suffices to prove that if \(v\) is a G-morphism from \(\pi\) to \(\varrho\) , then its image is contained in \(\mathrm{F}\) . Let \(p\) be the orthogonal projector of \(\mathrm{E}\) with image \(\mathrm{F}^{\circ}\) ; it is a G-morphism, since \(\mathrm{F}\) is a subrepresentation of \(\mathrm{E}\) (prop. 4 of V, p. 383). The image of the G-morphism \(p \circ v\) is orthogonal to \(\mathrm{F}\) ; it is therefore zero (otherwise \(\mathrm{C} \cup \{p \circ v\} \in \mathcal{O}\) , which contradicts the maximality of \(\mathrm{C}\) ), and consequently the image of \(v\) is contained in \(\mathrm{F}\) .

Now let \((\mathrm{E}_{i})_{i\in\mathrm{I}}\) and \((\mathrm{F}_{j})_{j\in\mathrm{J}}\) be families of pairwise orthogonal closed invariant subspaces of E, such that the subrepresentation of G in \(E_{i}\) (resp. in \(F_{j}\) ) is isomorphic to \(\pi\) for all \(i\in I\) (resp. for every \(j\in J\) ), and such that \(\mathrm{M}_{\pi}(\varrho)\) is the Hilbert sum of the family \((\mathrm{E}_{i})_{i\in\mathrm{I}}\) and of the family \((\mathrm{F}_{j})_{j\in\mathrm{J}}\) .

If I is finite, then

\[
\dim \operatorname{Hom} _ {\mathrm{G}} (\pi , \mathrm{M} _ {\pi} (\varrho)) = \dim \operatorname{Hom} _ {\mathrm{G}} \Bigl (\mathrm{E} _ {\pi}, \bigoplus_ {i \in \mathrm{I}} \mathrm{E} _ {i} \Bigr) = \operatorname{Card} (\mathrm{I})
\]

(formula (1) of V, p. 377 and cor. 2 of V, p. 387). For every finite subset L of J, we then have

\[
\begin{array}{c} \operatorname{Card} (\mathrm{L}) = \dim \operatorname{Hom} _ {\mathrm{G}} \Bigl (\mathrm{E} _ {\pi}, \bigoplus_ {j \in \mathrm{L}} \mathrm{F} _ {j} \Bigr) \\ \leqslant \dim \operatorname{Hom} _ {\mathrm{G}} (\pi , \operatorname{M} _ {\pi} (\varrho)) = \operatorname{Card} (\mathrm{I}) \end{array}
\]

(loc. cit.). This proves that J is then finite and that \(\text{Card(I)} = \text{Card(J)}\) , as desired. Likewise, if J is finite, then I is finite and has the same cardinal.

Now suppose that I and J are infinite. For \(j \in J\) , denote by \(p_{j}\) the orthogonal projector of E with image \(F_{j}\) . For \(i \in I\) , let \(x_{i}\) be a nonzero element of \(E_{i}\) ; it is a cyclic vector of \(E_{i}\) . Observe that since \(p_{j}\), induced by passage to the subspaces, is a G-morphism from \(E_{i}\) to \(F_{j}\) , the space \(p_{j}(E_{i})\) is zero if and only if \(p_{j}(x_{i}) = 0\) (indeed, the space \(E_{i} \cap \text{Ker}(p_{j})\) is either zero or equal to \(E_{i}\) ).

For every \(j \in J\) , there exists an element \(f(j) \in I\) such that \(p_{j}(\mathrm{E}_{f(j)})\) is not reduced to 0 (otherwise, the orthogonal projector \(p_{j}\) would be zero on \(\mathrm{M}_{\pi}(\varrho)\) ). We have thus defined a map f from J to I. For every \(i \in I\) , the set \(f^{-1}(i)\) is countable. Indeed, this set is contained in the set of \(j \in J\) such that \(p_{j}(\mathrm{E}_{i})\) is nonzero, that is, those such that \(p_{j}(x_{i}) \neq 0\) . Now, by (EVT, V, p. 20, cor. 2), we have

\[
\sum_ {j \in \mathrm{J}} \| p _ {j} (x _ {i}) \| ^ {2} = \| x _ {i} \| ^ {2},
\]

hence the set of \(j \in J\) such that \(p_{j}(x_{i}) \neq 0\) is countable. By E, III, p. 50, prop. 4, we conclude that \(\text{Card}(\text{J}) = \text{Card}(f(\text{J})) \leqslant \text{Card}(\text{I})\) . By interchanging the roles of I and J, we conclude that \(\text{Card}(\text{I}) = \text{Card}(\text{J})\) .

COROLLARY. --- Let G be a topological group and \(\varrho\) be a unitary representation of G in a complex Hilbert space E. The Hilbert sum of the \(\pi\) -isotypic components of G for \(\pi \in \widehat{\mathbf{G}}\) is the largest semisimple subrepresentation of \(\varrho\) .

Indeed, the \(\pi\) -isotypic components of G for \(\pi \in \widehat{\mathbf{G}}\) are pairwise orthogonal (prop. 8, \(a\) ), and each is semisimple (prop. 10, \(a\) ), hence the Hilbert sum F of the spaces \(\mathrm{M}_{\pi}(\varrho)\) for \(\pi \in \widehat{\mathbf{G}}\) defines a semisimple subrepresentation of \(\varrho\) . Since moreover every irreducible subrepresentation of \(\varrho\) is a subrepresentation of an isotypic component of \(\varrho\) , the corollary follows.

DEFINITION 10. --- Let G be a topological group. Let \(\varrho\) be a unitary representation of G in a complex Hilbert space E and \(\pi\) be an irreducible unitary representation of G in a complex Hilbert space.

One calls the multiplicity of \(\pi\) in \(\varrho\) the cardinal of the set I for any family \((\mathrm{E}_i)_{i\in \mathrm{I}}\) of closed subspaces of E stable under G such that the subrepresentation of \(\varrho\) induced in \(\mathrm{E}_i\) is isomorphic to \(\pi\) for every \(i\in I\) and such that \(\mathrm{M}_{\pi}(\varrho)\) is the Hilbert sum of the subspaces \(\mathrm{E}_i\) .

If the multiplicity of \(\pi\) in \(\varrho\) is finite, then it is equal to the dimension of the space \(\mathrm{Hom}_{\mathrm{G}}(\pi, \varrho)\) (resp. to the dimension of \(\mathrm{Hom}_{\mathrm{G}}(\varrho, \pi)\) ) by formula (1) of V, p. 377 and corollary 2 of V, p. 387. This is not always the case in general.

Remark. --- It is possible that a unitary representation \(\varrho\) is nonzero but that all isotypic components of \(\varrho\) relative to all irreducible representations of G are zero (cf. V, p. 426, remark).

